\documentclass[11pt,a4paper]{article}

\usepackage[utf8]{inputenc}
\usepackage[T1]{fontenc}
\usepackage[spanish]{babel}
\usepackage{amsmath,amssymb}
\usepackage{algorithm}
\usepackage{algpseudocode}
\usepackage{geometry}
\usepackage{hyperref}
\geometry{margin=2.5cm}

% Palabras clave del pseudocódigo en español.
\algrenewcommand\algorithmicrequire{\textbf{Entrada:}}
\algrenewcommand\algorithmicensure{\textbf{Salida:}}
\algrenewcommand\algorithmicreturn{\textbf{devolver}}
\algrenewcommand\algorithmicif{\textbf{si}}
\algrenewcommand\algorithmicthen{\textbf{entonces}}
\algrenewcommand\algorithmicelse{\textbf{si no}}
\algrenewcommand\algorithmicend{\textbf{fin}}
\algrenewcommand\algorithmicfor{\textbf{para}}
\algrenewcommand\algorithmicforall{\textbf{para cada}}
\algrenewcommand\algorithmicdo{\textbf{hacer}}
\algrenewcommand\algorithmicwhile{\textbf{mientras}}
\algrenewcommand\algorithmicfunction{\textbf{función}}
\algrenewcommand\algorithmicreturn{\textbf{devolver}}

\title{Random Forest básico para la clasificación de latidos ECG\\
\large Árbol de decisión, su entrenamiento y el bosque}
\author{Proyecto ECG Viewer --- paquete \texttt{ecg.ml}}
\date{}

\begin{document}
\maketitle

\begin{abstract}
Este documento describe el algoritmo básico del \emph{Random Forest} implementado a mano en
el paquete \texttt{ecg.ml} (sin librerías de \emph{machine learning}). Se explica, con
pseudocódigo, el árbol de decisión, su entrenamiento mediante entropía y corte por la media, y
el bosque como promedio de muchos árboles. El clasificador asigna a cada latido una de las
cuatro clases AAMI: \textbf{N} (normal), \textbf{S} (supraventricular), \textbf{V}
(ventricular) y \textbf{F} (fusión).
\end{abstract}

\section{Idea general}

Un \textbf{árbol de decisión} aprende una secuencia de preguntas del tipo
``¿la característica $j$ es menor que un umbral $u$?'' que separan los latidos por su clase. Cada
hoja del árbol guarda la proporción de clases de las muestras que cayeron en ella.

Un solo árbol tiende a \emph{sobreajustar} (memoriza el conjunto de entrenamiento). El
\textbf{Random Forest} (bosque aleatorio) combate esto entrenando muchos árboles distintos y
promediando sus predicciones (voto \emph{blando}). Los árboles se diferencian por dos fuentes de
azar:
\begin{enumerate}
  \item \textbf{Bootstrap}: cada árbol se entrena con una muestra de las filas tomada
        \emph{con reemplazo} (algunas filas se repiten, otras no aparecen).
  \item \textbf{Submuestreo de características}: en cada nodo solo se considera un subconjunto
        aleatorio de las características (típicamente $\sqrt{d}$ de las $d$ totales).
\end{enumerate}
Como los árboles cometen errores poco correlacionados, al promediarlos esos errores se
compensan y el bosque generaliza mejor que un árbol aislado.

\section{Entropía y ganancia de información}

La calidad de una partición se mide con la \textbf{entropía de Shannon}, que cuantifica el
``desorden'' de clases de un grupo de muestras $S$:
\[
  H(S) \;=\; -\sum_{c=1}^{K} p_c \,\log_2 p_c ,
  \qquad p_c = \frac{\#\{\text{muestras de clase } c\}}{|S|},
\]
donde $K$ es el número de clases ($K=4$ aquí) y los términos con $p_c = 0$ se omiten. Vale
$H(S)=0$ cuando todas las muestras son de la misma clase (grupo \emph{puro}) y es máxima cuando
están repartidas por igual.

Al partir $S$ en dos grupos $S_{\text{izq}}$ y $S_{\text{der}}$ según una pregunta, la
\textbf{ganancia de información} es la reducción de entropía:
\[
  \mathrm{GI}(S, j, u) \;=\; H(S)
  \;-\; \frac{|S_{\text{izq}}|}{|S|}\,H(S_{\text{izq}})
  \;-\; \frac{|S_{\text{der}}|}{|S|}\,H(S_{\text{der}}).
\]
La mejor pregunta de un nodo es la de mayor ganancia. En esta implementación el umbral $u$ de
una característica $j$ no se busca exhaustivamente: se fija como la \textbf{media} de esa
característica en las muestras del nodo (una elección simple y barata de calcular).

\begin{algorithm}[h]
\caption{Entropía de un grupo de muestras}
\begin{algorithmic}[1]
\Function{Entropia}{$S$, $K$}
  \Require conjunto de muestras $S$; número de clases $K$
  \Ensure entropía $H(S)$ en bits
  \If{$S$ está vacío} \Return $0$ \EndIf
  \State $\mathit{conteo}[1..K] \gets 0$
  \ForAll{muestra $m \in S$}
     \State $\mathit{conteo}[\,\text{clase}(m)\,] \mathrel{+}= 1$
  \EndFor
  \State $H \gets 0$;\quad $n \gets |S|$
  \For{$c \gets 1$ \textbf{hasta} $K$}
     \If{$\mathit{conteo}[c] > 0$}
        \State $p \gets \mathit{conteo}[c] / n$
        \State $H \gets H - p \cdot \log_2 p$
     \EndIf
  \EndFor
  \State \Return $H$
\EndFunction
\end{algorithmic}
\end{algorithm}

\section{El árbol de decisión}

Un nodo es \emph{interno} (una pregunta con dos hijos) o una \emph{hoja} (que guarda la
distribución de clases). La construcción es recursiva: se detiene y crea una hoja cuando el
grupo es pequeño (menos de $\ell$ muestras, el parámetro \texttt{min\_hoja}), cuando el grupo ya
es puro, o cuando ningún corte logra separar las muestras.

\subsection{Entrenamiento (construcción recursiva)}

\begin{algorithm}[h]
\caption{Construcción de un árbol de decisión}
\begin{algorithmic}[1]
\Function{Construir}{$S$}
  \Require muestras $S$; parámetros globales: $K$ clases, $m$ características por nodo,
           $\ell$ tamaño mínimo de hoja, generador aleatorio $\mathit{rng}$
  \Ensure un nodo (raíz del subárbol)
  \If{$|S| < \ell$ \textbf{o} \Call{MismaClase}{$S$}}
     \State \Return hoja con distribución \Call{Distribucion}{$S$}
  \EndIf
  \State $d \gets$ número de características de una muestra
  \State $C \gets$ \Call{MuestraSinReemplazo}{$d, m, \mathit{rng}$}
        \Comment{$m$ características candidatas al azar}
  \State $H_{\text{padre}} \gets$ \Call{Entropia}{$S, K$}
  \State $(j^\ast, u^\ast, g^\ast) \gets (\text{indefinido}, 0, -\infty)$
  \ForAll{característica $j \in C$}
     \State $u \gets$ media de la característica $j$ en $S$
     \State $g \gets$ \Call{GananciaInformacion}{$S, j, u, H_{\text{padre}}$}
     \If{$g > g^\ast$}
        \State $(j^\ast, u^\ast, g^\ast) \gets (j, u, g)$
     \EndIf
  \EndFor
  \State $(S_{\text{izq}}, S_{\text{der}}) \gets$ \Call{Repartir}{$S, j^\ast, u^\ast$}
  \If{$S_{\text{izq}}$ vacío \textbf{o} $S_{\text{der}}$ vacío}
     \State \Return hoja con distribución \Call{Distribucion}{$S$}
        \Comment{el corte no separa nada}
  \EndIf
  \State $\mathit{nodo}.\text{caracteristica} \gets j^\ast$;\quad
         $\mathit{nodo}.\text{umbral} \gets u^\ast$
  \State $\mathit{nodo}.\text{izquierda} \gets$ \Call{Construir}{$S_{\text{izq}}$}
  \State $\mathit{nodo}.\text{derecha} \gets$ \Call{Construir}{$S_{\text{der}}$}
  \State \Return $\mathit{nodo}$
\EndFunction
\end{algorithmic}
\end{algorithm}

\begin{algorithm}[h]
\caption{Reparto de muestras y ganancia de información}
\begin{algorithmic}[1]
\Function{Repartir}{$S, j, u$}
  \State $S_{\text{izq}} \gets \emptyset$;\quad $S_{\text{der}} \gets \emptyset$
  \ForAll{muestra $m \in S$}
     \If{$m.\text{caracteristica}[j] < u$}
        \State añadir $m$ a $S_{\text{izq}}$
     \Else
        \State añadir $m$ a $S_{\text{der}}$
     \EndIf
  \EndFor
  \State \Return $(S_{\text{izq}}, S_{\text{der}})$
\EndFunction
\Statex
\Function{GananciaInformacion}{$S, j, u, H_{\text{padre}}$}
  \State $(S_{\text{izq}}, S_{\text{der}}) \gets$ \Call{Repartir}{$S, j, u$}
  \State $n \gets |S|$
  \State \Return $H_{\text{padre}}
        - \dfrac{|S_{\text{izq}}|}{n}\,\Call{Entropia}{S_{\text{izq}}, K}
        - \dfrac{|S_{\text{der}}|}{n}\,\Call{Entropia}{S_{\text{der}}, K}$
\EndFunction
\Statex
\Function{Distribucion}{$S$}
  \Comment{proporción de cada clase: lo que guarda una hoja}
  \State $\mathit{dist}[1..K] \gets 0$
  \ForAll{muestra $m \in S$}
     \State $\mathit{dist}[\,\text{clase}(m)\,] \mathrel{+}= 1 / |S|$
  \EndFor
  \State \Return $\mathit{dist}$
\EndFunction
\end{algorithmic}
\end{algorithm}

\subsection{Predicción con un árbol}

Para clasificar un vector de características $x$ se baja por el árbol desde la raíz: en cada
nodo interno se va a la izquierda si $x[\,\text{caracteristica}\,] < \text{umbral}$ y a la
derecha en caso contrario, hasta llegar a una hoja. Se devuelve la distribución de clases de
esa hoja.

\begin{algorithm}[h]
\caption{Predicción de un árbol}
\begin{algorithmic}[1]
\Function{PredecirArbol}{$x$}
  \Require vector de características $x$
  \Ensure distribución de clases $\mathit{dist}[1..K]$
  \State $\mathit{nodo} \gets \text{raíz}$
  \While{$\mathit{nodo}$ no es hoja}
     \If{$x[\,\mathit{nodo}.\text{caracteristica}\,] < \mathit{nodo}.\text{umbral}$}
        \State $\mathit{nodo} \gets \mathit{nodo}.\text{izquierda}$
     \Else
        \State $\mathit{nodo} \gets \mathit{nodo}.\text{derecha}$
     \EndIf
  \EndWhile
  \State \Return $\mathit{nodo}.\text{distribucion}$
\EndFunction
\end{algorithmic}
\end{algorithm}

\section{El bosque (Random Forest)}

El bosque guarda una lista de árboles. Para \textbf{entrenar}, por cada árbol se toma un
bootstrap de las muestras y se construye un árbol con él (todos los árboles comparten el mismo
generador aleatorio $\mathit{rng}$, de modo que con la misma semilla se obtiene el mismo
bosque). Para \textbf{predecir}, se promedian las distribuciones de todos los árboles y se
devuelve la clase de mayor probabilidad media (empate $\to$ la clase de menor índice).

\begin{algorithm}[h]
\caption{Entrenamiento del bosque}
\begin{algorithmic}[1]
\Function{EntrenarBosque}{$S$}
  \Require muestras de entrenamiento $S$; parámetros: $T$ árboles, $K$ clases,
           $m$ características por nodo, $\ell$ mínimo por hoja, semilla $s$
  \Ensure bosque entrenado (lista de $T$ árboles)
  \State $\mathit{rng} \gets$ generador aleatorio sembrado con $s$
  \State $\text{arboles} \gets$ lista vacía
  \For{$t \gets 1$ \textbf{hasta} $T$}
     \State $I \gets$ \Call{Bootstrap}{$|S|, \mathit{rng}$}
        \Comment{$|S|$ índices al azar con reemplazo}
     \State $B \gets \{\, S[i] : i \in I \,\}$
     \State $\text{arbol} \gets$ \Call{Construir}{$B$}
        \Comment{usa $K, m, \ell, \mathit{rng}$}
     \State añadir $\text{arbol}$ a $\text{arboles}$
  \EndFor
  \State \Return $\text{arboles}$
\EndFunction
\end{algorithmic}
\end{algorithm}

\begin{algorithm}[h]
\caption{Predicción del bosque (voto blando)}
\begin{algorithmic}[1]
\Function{PredecirDistribucion}{$x$}
  \State $\mathit{acum}[1..K] \gets 0$
  \ForAll{$\text{arbol} \in \text{arboles}$}
     \State $\mathit{dist} \gets$ \Call{PredecirArbol}{$\text{arbol}, x$}
     \For{$c \gets 1$ \textbf{hasta} $K$}
        \State $\mathit{acum}[c] \mathrel{+}= \mathit{dist}[c]$
     \EndFor
  \EndFor
  \For{$c \gets 1$ \textbf{hasta} $K$}
     \State $\mathit{acum}[c] \gets \mathit{acum}[c] / |\text{arboles}|$
  \EndFor
  \State \Return $\mathit{acum}$
\EndFunction
\Statex
\Function{Predecir}{$x$}
  \Ensure índice de clase predicha ($1 \ldots K$)
  \State $\mathit{dist} \gets$ \Call{PredecirDistribucion}{$x$}
  \State \Return $\arg\max_c \; \mathit{dist}[c]$
\EndFunction
\end{algorithmic}
\end{algorithm}

\subsection{Las dos fuentes de azar}

\begin{algorithm}[h]
\caption{Muestreo aleatorio del bosque}
\begin{algorithmic}[1]
\Function{Bootstrap}{$n, \mathit{rng}$}
  \Comment{$n$ índices en $[0,n)$ al azar \textbf{con} reemplazo}
  \State $I[0..n-1]$
  \For{$i \gets 0$ \textbf{hasta} $n-1$}
     \State $I[i] \gets \mathit{rng}.\text{entero}(0, n)$
  \EndFor
  \State \Return $I$
\EndFunction
\Statex
\Function{MuestraSinReemplazo}{$\mathit{total}, k, \mathit{rng}$}
  \Comment{$k$ índices \textbf{distintos} de $[0,\mathit{total})$: Fisher--Yates parcial}
  \State $k \gets \min(k, \mathit{total})$
  \State $a[i] \gets i$ para $i \in [0, \mathit{total})$
  \For{$i \gets 0$ \textbf{hasta} $k-1$}
     \State $j \gets i + \mathit{rng}.\text{entero}(0, \mathit{total}-i)$
     \State intercambiar $a[i] \leftrightarrow a[j]$
  \EndFor
  \State \Return $a[0 \ldots k-1]$
\EndFunction
\end{algorithmic}
\end{algorithm}

\section{Complejidad y parámetros}

Sea $n = |S|$ el número de latidos de entrenamiento, $d$ el número de características
(aquí $d = 257$: $252$ de forma de onda más $5$ de ritmo RR), $T$ el número de árboles y
$m \approx \sqrt{d}$ las características consideradas por nodo.

\begin{itemize}
  \item \textbf{Construir un nodo} evalúa $m$ cortes y, por cada uno, reparte las muestras y
        calcula dos entropías: coste $O(m \cdot n_{\text{nodo}})$ sobre las $n_{\text{nodo}}$
        muestras del nodo.
  \item \textbf{Un árbol} tiene profundidad $\approx \log(n/\ell)$; sumando los niveles, el
        coste es del orden de $O(m \, n \log n)$.
  \item \textbf{El bosque} multiplica por $T$: $O(T \, m \, n \log n)$.
  \item \textbf{Predecir} un latido con el bosque es $O(T \cdot \text{profundidad})$.
\end{itemize}

Valores típicos usados en la aplicación: $T$ entre $15$ y $50$ árboles,
$m = \lfloor \sqrt{257} \rfloor = 16$ características por nodo, $\ell$ (mínimo por hoja) entre
$10$ y $20$. Con las $\sim\!51\,000$ latidos de la partición DS1 de MIT-BIH, el entrenamiento
tarda del orden de decenas de segundos; por eso se ejecuta en segundo plano y el modelo
entrenado se puede \emph{guardar en disco} (formato de texto \texttt{RFV1}: un volcado en
preorden de cada árbol) para reutilizarlo entre sesiones sin re-entrenar.

\section{Correspondencia con el código}

\begin{tabular}{@{}ll@{}}
\textbf{Pseudocódigo} & \textbf{Clase / método Java} \\
\hline
\textsc{Entropia} & \texttt{ecg.ml.Entropia.entropia} \\
\textsc{Construir}, \textsc{Repartir}, \textsc{GananciaInformacion} & \texttt{ecg.ml.ArbolDesicion} \\
\textsc{PredecirArbol} & \texttt{ArbolDesicion.predecir} \\
\textsc{EntrenarBosque} & \texttt{ecg.ml.BosqueSimple.entrenar} \\
\textsc{PredecirDistribucion}, \textsc{Predecir} & \texttt{BosqueSimple.predecirDistribucion / predecir} \\
\textsc{Bootstrap}, \textsc{MuestraSinReemplazo} & \texttt{ecg.ml.Muestreo} \\
Guardar / cargar modelo & \texttt{BosqueSimple.guardar / cargar} \\
\end{tabular}

\end{document}
